\section{Setup and assumptions}

Throughout, \leanref{sym:n}{$n$} denotes sample size and \leanref{sym:d}{$d$} denotes the number of adjustment cells. Cell indices such as \leanref{sym:k}{$k$} range over \(\{1,\ldots,d\}\), and treatment levels such as \leanref{sym:a}{$a$} range over \(\{0,1\}\). Constants \leanref{sym:c_epsilon}{$c_\epsilon$} and \leanref{sym:C_epsilon}{$C_\epsilon$} may change from line to line and depend only on the overlap parameter \leanref{sym:epsilon}{$\epsilon$}. The outcome scale is \leanref{sym:M}{$M$}\(\ge 1\), and the heterogeneity radius is \leanref{sym:sigma}{$\sigma$}\(\in[0,2]\). We use \(O(\cdot)\), \(\lesssim_\epsilon\), and \(\asymp_\epsilon\) with constants depending only on \(\epsilon\).

A full-data law \leanref{sym:P}{$P$} governs the finite-cell observational unit with confounder \leanref{sym:X}{$X$}, binary treatment \leanref{sym:A}{$A$}, potential outcomes \leanref{sym:Y(a)}{$Y(a)$}, and observed real outcome \leanref{sym:Y}{$Y$}. For a supported cell, \leanref{sym:p_k}{$p_k$} is the cell mass, \leanref{sym:pi_k}{$\pi_k$} is the propensity, and \leanref{sym:mu_{ak}}{$\mu_{ak}$} is the arm-cell conditional mean. The corresponding cell treatment effect is \leanref{sym:tau_k}{$\tau_k$}, the target average treatment effect is \leanref{sym:tau(P)}{$\tau(P)$}, and the centered cell-effect deviation is \leanref{sym:delta_k}{$\delta_k$}. These objects are fixed formally below before the minimax experiment is introduced.

The analysis starts from the potential-outcomes conditions for finite-cell adjustment, followed by scale and heterogeneity restrictions that make the finite-sample risk problem comparable across real-outcome laws.

\begin{assumptionv}{ass:consistency}[Consistency]
Under \(P\),
\[
Y=Y(A)
\]
almost surely.
\end{assumptionv}

Consistency links the observed outcome to the potential outcome under the realized treatment assignment, following the potential-outcomes setup used in discrete-adjustment analyses such as \citet{ZengBalakrishnanHanKennedy2024}.

\begin{assumptionv}{ass:conditional-exchangeability}[Conditional exchangeability]
Under \(P\),
\[
(Y(0),Y(1))\perp A\mid X .
\]
\end{assumptionv}

Conditional exchangeability is the selection-on-observables condition after conditioning on the discrete adjustment cell \citep{ZengBalakrishnanHanKennedy2024}. Together with consistency, it gives the usual causal interpretation of cell-level treated-control contrasts \citep{Rubin1974,RosenbaumRubin1983,ImbensRubin2015}.

\begin{assumptionv}{ass:overlap}[Fixed overlap]
For each cell \(k\), let
\[
p_k:=P(X=k),
\]
and, when \(p_k>0\), let
\[
\pi_k:=P(A=1\mid X=k).
\]
For every \(k\) with \(p_k>0\),
\[
\epsilon\le \pi_k\le 1-\epsilon .
\]
\end{assumptionv}

The fixed-overlap requirement imposes strong overlap on every supported cell, as in binary discrete-adjustment work such as \citet{ZengBalakrishnanHanKennedy2024}. It keeps both treatment arms available in population within each cell while allowing the cell masses themselves to be arbitrary.

\begin{assumptionv}{ass:mean-normalization}[Mean normalization]
The law \(P\) carries an arm-cell outcome distribution for every arm \(a\in\{0,1\}\) and \emph{every} cell \(k\), zero-mass cells included; let \(\mu_{ak}\) denote its mean.  Whenever the arm-cell probability \(P(A=a,X=k)\) is positive, \(\mu_{ak}\) is the observed conditional mean \(E_P[Y\mid A=a,X=k]\).  Only supported cells are constrained: for every \(a\in\{0,1\}\) and every \(k\) with \(p_k>0\),
\[
|\mu_{ak}|\le M/2 .
\]
\end{assumptionv}

Mean normalization is the affine-centered bounded conditional-mean analogue of the bounded binary response scale in \citet{ZengBalakrishnanHanKennedy2024}. The centering fixes the outcome scale for the minimax comparison, and the restriction applies on the supported part of the adjustment distribution.

\begin{assumptionv}{ass:approximate-homogeneity}[Approximate homogeneity]
For \emph{every} cell \(k\), let
\[
\tau_k:=\mu_{1k}-\mu_{0k}
\]
and
\[
\delta_k:=\tau_k-\sum_j p_j\tau_j ,
\]
both well defined on zero-mass cells as well.
The cell-effect deviations satisfy
\[
\max_{k:p_k>0}|\delta_k|\le \sigma M .
\]
\end{assumptionv}

Approximate homogeneity is the outcome-scale-normalized analogue of a maximal cell-effect-deviation radius in binary discrete adjustment \citep{ZengBalakrishnanHanKennedy2024}. The radius \(\sigma M\) measures the largest supported-cell departure of a cell treatment effect from the population-weighted average effect.

\begin{assumptionv}{ass:second-central-moment}[Second central moment]
For every \(a\in\{0,1\}\) and every \(k\) with \(p_k>0\),
\[
E_P\!\left[(Y-\mu_{ak})^2\mid A=a,X=k\right]\le M^2 .
\]
\end{assumptionv}

The moment bound extends the binary bounded-outcome scale to real-valued outcomes through a uniform conditional second-central-moment envelope. It keeps the noise level on the same scale as the conditional-mean and heterogeneity restrictions.

These conditions define two law classes. The first imposes the radius bound, while the second records the same overlap and moment envelope with the radius left free.

\begin{definitionv}{def:model-class}[Real-outcome model class \(\mathcal P_{d,\epsilon,M,\sigma}\)]
\[
\mathcal P_{d,\epsilon,M,\sigma}
:=
\left\{
P:
\begin{array}{l}
X\in\{1,\ldots,d\},\ A\in\{0,1\},\ Y,Y(0),Y(1)\in\mathbb R,\\
0<\epsilon<1/2,\quad M\ge 1,\quad 0\le \sigma\le 2,\\
P\text{ satisfies }\cref{obj:ass:consistency,obj:ass:conditional-exchangeability,obj:ass:overlap,obj:ass:mean-normalization,obj:ass:second-central-moment,obj:ass:approximate-homogeneity}
\end{array}
\right\}.
\]
\end{definitionv}

The radius-indexed class \leanref{sym:\mathcal P_{d,epsilon,M,sigma}}{$\mathcal P_{d,\epsilon,M,\sigma}$} is the main parameter space of the paper: all finite-cell masses are allowed, and the scalar radius records the amount of cell-effect variation available to the estimator.

\begin{definitionv}{def:unrestricted-class}[Unrestricted-radius class \(\mathcal P_{d,\epsilon,M}^{\mathrm{unr}}\)]
\[
\mathcal P_{d,\epsilon,M}^{\mathrm{unr}}
:=
\left\{
P:
\begin{array}{l}
X\in\{1,\ldots,d\},\ A\in\{0,1\},\ Y,Y(0),Y(1)\in\mathbb R,\\
0<\epsilon<1/2,\quad M\ge 1,\\
P\text{ satisfies }\cref{obj:ass:consistency,obj:ass:conditional-exchangeability,obj:ass:overlap,obj:ass:mean-normalization,obj:ass:second-central-moment}
\end{array}
\right\}.
\]
\end{definitionv}

The unrestricted-radius class \leanref{sym:\mathcal P_{d,epsilon,M}^{unr}}{$\mathcal P_{d,\epsilon,M}^{\mathrm{unr}}$} supplies the endpoint comparison for the radius scale under the same fixed-overlap and real-outcome moment envelope.

The observed-data experiment is generated by independent sampling from the observed margin. Write the observed record as \leanref{sym:O}{$O$} and its observed-data law as \leanref{sym:P_O}{$P_O$}.

\begin{definitionv}{def:sample-experiment}[Sample experiment \(\mathcal E_{n,d,\epsilon,M,\sigma}\)]
Let
\[
O:=(X,A,Y)\in\mathcal O:=\{1,\ldots,d\}\times\{0,1\}\times\mathbb R,
\]
and let \(P_O\) denote the observed-data law of \(O\) under \(P\). The \(n\)-sample observed-data experiments generated by \(\mathcal P_{d,\epsilon,M,\sigma}\) are
\[
\mathcal E_{n,d,\epsilon,M,\sigma}
:=
\left\{
Q_P^{(n)}=P_O^{\otimes n}:P\in\mathcal P_{d,\epsilon,M,\sigma}
\right\}.
\]
\end{definitionv}

Thus each sampling law \leanref{sym:Q_P^{(n)}}{$Q_P^{(n)}$} belongs to the experiment \leanref{sym:\mathcal E_{n,d,epsilon,M,sigma}}{$\mathcal E_{n,d,\epsilon,M,\sigma}$} induced by a law in \(\mathcal P_{d,\epsilon,M,\sigma}\).

The target parameter is the finite-cell average treatment effect. Its definition uses the same cell masses and arm-cell conditional means introduced in the assumptions.

\begin{definitionv}{def:ate-functional}[Average treatment effect \(\tau(P)\)]
For a real-outcome finite-cell law \(P\) indexed by \(k=1,\ldots,d\) and an overlap parameter \(\epsilon\) with \(0<\epsilon<1/2\), assume consistency, conditional exchangeability of \((Y(0),Y(1))\) and \(A\) given \(X\) for measurable potential-outcome events, overlap \(\epsilon\le \pi_k\le 1-\epsilon\) whenever \(p_k>0\), and first-moment integrability of each arm-cell outcome law for every \(a\in\{0,1\}\) and positive-mass cell. With \(p_k\) the cell mass and \(\mu_{ak}\) the arm-cell conditional outcome mean, define
\[
\tau(P)
:=
\sum_{k=1}^{d}p_k(\mu_{1k}-\mu_{0k}).
\]
\end{definitionv}

\begin{propositionv}{prop:ate-identification}[ATE g-formula identification]
For every dimension \(d\), overlap parameter \(\epsilon\), envelope level \(M\), radius \(\sigma\), and every \(P\in\mathcal P_{d,\epsilon,M,\sigma}\) in the sense of \cref{obj:def:model-class}, the average treatment effect functional \(\tau(P)\) of \cref{obj:def:ate-functional} equals the full-data causal contrast and is identified from the observed margin by the finite-cell g-formula:
\[
\tau(P)
=
\int \{Y(1)-Y(0)\}\,dP
=
\sum_{k=1}^{d} p_k
\left\{
E_P[Y\mid A=1,X=k]-E_P[Y\mid A=0,X=k]
\right\}.
\]
\end{propositionv}

\Cref{obj:prop:ate-identification} separates causal identification from the finite-sample estimation problem. Under the stated conditions, the minimax analysis can target the observed-data g-formula representation of the full-data causal contrast, in the tradition of adjustment and g-computation arguments \citep{Rubin1974,RosenbaumRubin1983,Robins1994,ImbensRubin2015}.

For the risk statements, estimators are clipped to the natural scale implied by the mean normalization. The selector benchmark used later is
\[
r_{n,d,\sigma}
:=
n^{-1}
+
\min\left\{
\sigma^2+\frac{d}{n^2},
\frac{d^2}{n^2\log^2(en)},
1
\right\},
\]
where \leanref{sym:r_{n,d,sigma}}{$r_{n,d,\sigma}$} records the parametric term together with the best of the collision, polynomial, and zero-estimator scales.

\begin{definitionv}{def:minimax-risk}[Minimax risk \(\mathsf R_{n,d,\epsilon,M,\sigma}\) and estimator class \(\mathcal T_{n,M}\)]
Let
\[
\mathcal T_{n,M}
:=
\bigl\{
T:\mathcal O^n\to[-M,M]\;:\;T\ \text{is measurable}
\bigr\}.
\]
The minimax mean-squared risk over \(\mathcal P_{d,\epsilon,M,\sigma}\) is
\[
\mathsf R_{n,d,\epsilon,M,\sigma}
:=
\inf_{T\in\mathcal T_{n,M}}
\sup_{P\in\mathcal P_{d,\epsilon,M,\sigma}}
E_{Q_P^{(n)}}\!\left[(T-\tau(P))^2\right].
\]
\end{definitionv}

The estimator class \leanref{sym:\mathcal T_{n,M}}{$\mathcal T_{n,M}$} matches the range of the target under the scale bound, and the minimax risk \leanref{sym:\mathsf R_{n,d,epsilon,M,sigma}}{$\mathsf R_{n,d,\epsilon,M,\sigma}$} is the paper's central finite-sample loss criterion.

The main rate displays use the following benchmark notation. The table gives each symbol one role; later sections keep these names fixed.
\[
\begin{array}{lll}
\hline
\text{Symbol} & \text{Name} & \text{Definition}\\
\hline
u_{n,d} & \text{rare-cell polynomial remainder} & d^2/\{n^2\log^2(en)\}\\
h_{n,d,\sigma} & \text{collision remainder} & \sigma^2+d/n^2\\
b_{n,d} & \text{exact-homogeneity baseline} & n^{-1}+d/n^2\\
r_{n,d,\sigma} & \text{selector benchmark} & n^{-1}+\min\{1,u_{n,d},h_{n,d,\sigma}\}\\
q_{n,d,\sigma} & \text{triangular selector benchmark} & n^{-1}+\min\{u_{n,d},\sigma^2+d/n^2\}\\
\ell_{n,d,\sigma} & \text{product converse benchmark} & b_{n,d}+\sigma^2u_{n,d}\\
\hline
\end{array}
\]

The normalization also pins down the unrestricted endpoint of the radius scale.

\begin{lemmav}{lem:scale-sanity}[Scale sanity]
For any \(d\in\mathbb N\) and any \(\epsilon,M\in\mathbb R\), let \(\mathcal P_{d,\epsilon,M}^{\mathrm{unr}}\) and \(\mathcal P_{d,\epsilon,M,\sigma}\) be the model classes in \cref{obj:def:unrestricted-class,obj:def:model-class}. For each law \(P\), write
\[
\tau_k(P)=\mu_{1k}(P)-\mu_{0k}(P),
\qquad
\tau(P)=\sum_{k\in\{1,\ldots,d\}} p_k(P)\tau_k(P),
\qquad
\delta_k(P)=\tau_k(P)-\tau(P).
\]
Every \(P\in\mathcal P_{d,\epsilon,M}^{\mathrm{unr}}\) satisfies
\[
|\tau_k(P)|\le M\quad\text{for every }k\text{ with }p_k(P)>0,
\qquad
|\tau(P)|\le M,
\]
and
\[
|\delta_k(P)|\le 2M\quad\text{for every }k\text{ with }p_k(P)>0.
\]
Consequently the two classes have the same underlying laws at radius \(2\): for every \(P\in\mathcal P_{d,\epsilon,M}^{\mathrm{unr}}\) there exists \(Q\in\mathcal P_{d,\epsilon,M,2}\) with the same law as \(P\), and for every \(Q\in\mathcal P_{d,\epsilon,M,2}\) there exists \(P\in\mathcal P_{d,\epsilon,M}^{\mathrm{unr}}\) with the same law as \(Q\).
\end{lemmav}

\Cref{obj:lem:scale-sanity} makes the radius normalization exact: the bounded conditional means imply \(|\tau(P)|\le M\), supported cell effects lie within the same outcome scale, and radius \(2\) coincides with the unrestricted-radius law class. This calibration is used when the main bracket is specialized to endpoint regimes.

Finally, the lower-bound arguments later use binary source experiments from \citet{ZengBalakrishnanHanKennedy2024}. We record the source target and source law classes here so that the subsequent affine embeddings have a fixed reference point.

\begin{definitionv}{synth_1}[Binary exact-homogeneity source target]
For a binary-outcome fixed-sample law \(P\) in the uniform-mass exact-homogeneity source family, let \(\mu_{ak}\) denote the conditional response mean in arm \(a\in\{0,1\}\) and cell \(k\in\{1,\ldots,d\}\). The source target is\[\psi(P):=\frac{1}{d}\sum_{k=1}^{d}(\mu_{1k}-\mu_{0k}).\]
\end{definitionv}

\begin{definitionv}{synth_2}[Binary source law classes]
For an overlap parameter \(0<\epsilon<1/2\), let \(\mathcal D_d^{\mathrm{bin}}(\epsilon)\) be the class of binary \(d\)-cell full-data laws for \((X,A,B(0),B(1),B)\), with \(X\in\{1,\ldots,d\}\), \(A,B(0),B(1),B\in\{0,1\}\), consistency \(B=B(A)\), conditional exchangeability \((B(0),B(1))\perp A\mid X\), arbitrary cell masses \(p_k\), propensities \(\pi_k\) satisfying \(\epsilon\le\pi_k\le1-\epsilon\) on every positive-mass cell, and arm means \(\mu_{ak}^{\mathrm{bin}}=E[B\mid A=a,X=k]\in[0,1]\). Define its uniform-mass exactly homogeneous subclass by
\[
\mathcal H_d^{\mathrm{bin}}(\epsilon)
:=
\left\{
P\in\mathcal D_d^{\mathrm{bin}}(\epsilon):
p_k(P)=\frac1d\ \text{for every }k,\ 
\mu_{1k}^{\mathrm{bin}}(P)-\mu_{0k}^{\mathrm{bin}}(P)
=
\mu_{1\ell}^{\mathrm{bin}}(P)-\mu_{0\ell}^{\mathrm{bin}}(P)
\ \text{for all }k,\ell
\right\}.
\]
\end{definitionv}

The source target in \cref{obj:synth_1} and the binary classes in \cref{obj:synth_2} provide the binary ingredients transported in the lower-bound section. Their role is to specify the published discrete-adjustment experiments before the paper embeds them into the real-outcome class \(\mathcal P_{d,\epsilon,M,\sigma}\).
